\documentclass[10pt,journal]{IEEEtran}
\usepackage[T1]{fontenc}
\usepackage{amsmath,amssymb,booktabs,graphicx,url,cite}
\usepackage[colorlinks=true,linkcolor=blue,citecolor=blue,urlcolor=blue]{hyperref}
\title{Evaluating Model-Free Object Tracking and Surface Reconstruction for RGB-D Teleoperation: A Development Study}
\author{Tiwater Research (\href{mailto:research@tiwater.com}{\texttt{research@tiwater.com}})}
\begin{document}
\maketitle
\begin{abstract}
Persistent object pose and surface geometry are important for RGB-D teleoperation, but tracking availability does not establish accurate geometry or end-to-end readiness. We examine the existing BundleTrack-based pose owner in EvidenceScene and Point2Pose as an offline alternative. On six selected public YCBInEOAT sequences, we report failure-inclusive pose recall with per-frame reference masks, and separately investigate stored object-surface coverage, depth filtering, and repeatability. At a post-hoc 10 mm nearest-vertex threshold, BundleTrack achieves more passing frames than Point2Pose on each selected sequence. On sugar, the direction reverses under corresponding-vertex ADD: Point2Pose/BundleTrack pass 206/174 frames, versus 617/794 under ADD-S. On the additional bleach-hard sequence, the 5 mm count and 0--100 mm AUC favor Point2Pose. Thresholds and repeat runs also change observed differences. A public sugar-sequence depth-filter intervention produces 417--737 passes across two no-filter runs, compared with 729--794 across two runs using the existing filter, out of 1,001 post-initial frames. These selected reference-assisted results do not establish a general method ranking or an improvement from changing the filter. We identify the independent physical reference, held-out recordings, complete latency measurement, and release provenance needed for a replacement decision.
\end{abstract}
\begin{IEEEkeywords}RGB-D tracking, object reconstruction, reproducibility, teleoperation, failure-inclusive evaluation\end{IEEEkeywords}
\section{Introduction}
Teleoperation needs both object geometry and a reliable account of when that geometry is current. BundleTrack~\cite{bundletrack} tracks unfamiliar objects using RGB-D without an instance CAD model. Point2Pose~\cite{point2pose} combines point tracking with reconstruction and targets multi-object tracking and recovery after complete occlusion. Our selected single-object comparisons do not test its authors' recovery claims.

BundleSDF~\cite{bundlesdf} jointly tracks and reconstructs an unknown object from RGB-D and an initial mask. Point2Pose's authors report higher mean YCBInEOAT ADD and ADD-S AUC for BundleSDF, with one BundleSDF sequence marked failed. Our BundleTrack and causal surface are a different system; we also report a conditional BundleSDF development replay. FoundationPose~\cite{foundationpose} supports CAD or reference views; our diagnostic uses CAD.

EvidenceScene uses SAM~2, BundleTrack and a causal surface owner. Replacement requires shadow evaluation and a paired gate. We examine failure-inclusive tracking and persistent surfaces; we do not claim a new state-of-the-art method.

\section{Protocol}
\subsection{Pose comparison}
We analyze five selected YCBInEOAT~\cite{tracknet} object recordings: cracker box, mustard, bleach, sugar and tomato, plus a subsequently selected bleach-hard sequence. The Point2Pose implementation is pinned at commit \texttt{51856226610d}; it is not a benchmark of an unspecified current or future release. Both offline trackers receive the same decoded RGB-D and the dataset reference mask on every frame, with a common 3-by-3 erosion. Thus the study is reference-mask assisted and does not measure autonomous segmentation. Normalized boundary receipts are checked for the public runs; method-specific downstream preprocessing differs. BundleTrack's tested full-audit configuration differs from the production configuration. Runs were separate, and complete matched latency was not recorded.

Each tracker uses its own initial object gauge, aligned after inference using the first annotated pose; that frame is excluded. Later frames use reference-to-prediction nearest-vertex ADD-S, with loss or non-candidate outputs counted as failures. The 10 and 20 mm thresholds and normalized exact recall AUC over 0--100 mm were selected during development, not taken from an official or deployment protocol. We also compute corresponding-vertex ADD. Model vertices are nonuniform, symmetry groups are not explicitly enumerated, and annotations and masks are not independent physical truth.

\subsection{Surface and preprocessing studies}
We compare three snapshots on each of three inspected public sequences: Point2Pose TSDF meshes versus BundleTrack sparse voxel surfaces. Initial gauges use the first annotation without ICP. A sensitivity analysis maps both to a 4 mm grid and evaluates bidirectional nearest distances at 10 mm; representation, sampling and incomplete observations limit equivalence.

We privately replay the tracker on 552 bimanual chest frames with saved SAM~2 masks, testing its 95th-percentile depth cutoff against no filtering, twice per setting. A separate sugar experiment compares two runs per setting against annotations. These stateful interventions leave the live owner unchanged; tracker entry and filtered depth were instrumented on the first ten chest frames only.

\section{Results}
\subsection{Failure-inclusive pose recall}
Table~\ref{tab:pose} reports failure-inclusive counts for each complete post-initial denominator. At 10 mm the existing BundleTrack replay has more passing frames in each selected sequence. These are temporally correlated observations, not independent trials. On mustard, a post-hoc Point2Pose camera-local rotation jump of $34.85^\circ$ occurs between adjacent selected frames with a 200 ms RGB-header gap; BundleTrack and the annotations change $0.086^\circ$ and $0.036^\circ$. The association and annotations are not independent physical truth. At frames 707 and 710, Point2Pose annotation-gauge rotations differ by $28.67^\circ$ and $20.99^\circ$, yet ADD-S is 6.79 and 7.14 mm. The corresponding directed model-long-axis differences are $6.71^\circ$ and $6.82^\circ$; symmetry and task relevance are not established. For sugar, 187 BundleTrack-only passes form 88 contiguous runs, the longest 29 frames; runs themselves need not be independent. For tomato the paired differences are two Point2Pose-only and four BundleTrack-only frames.

The comparison depends on the post-hoc threshold and metric. On sugar at 5 mm Point2Pose/BundleTrack pass 194/167 frames, whereas at 10 mm they pass 617/794. On tomato at 6 mm they pass 1,100/1,065, compared with 1,301/1,303 at 10 mm. Cracker AUC also differs in direction from its 10 mm count comparison. On sugar, corresponding-vertex ADD at 10 mm passes 206/174 frames for Point2Pose/BundleTrack, compared with 617/794 for ADD-S. ADD penalizes model-coordinate rotation, including possible symmetries; neither metric alone establishes task accuracy.

Reported availability alone is not reference agreement. On bleach, Point2Pose and BundleTrack report 660/662 and 655/662 available outputs, yet 505 and 354 of those outputs respectively exceed the post-hoc 10 mm ADD-S criterion; at 20 mm the corresponding counts are 476 and 188. These status conventions differ between methods, and the reference is not independent physical truth.

On a subsequently selected bleach-hard recording, Point2Pose/BundleTrack pass 283/49 frames at 5 mm, 416/418 at 10 mm and 439/436 at 20 mm, out of 440 post-initial frames; normalized AUC is 0.94788/0.93035. At 10 mm, 9 frames pass only for Point2Pose and 11 only for BundleTrack; at 5 mm the respective counts are 242 and 8. Point2Pose reports no losses and BundleTrack two non-candidates. These are correlated, reference-mask-assisted development outcomes. Scoring follows the upstream sorted-file convention: only 2/441 pose filenames equal the bag RGB ordinal or extracted-frame index, and the maximum selected image-header gap is 5 seconds. The 441 pose files have 159 distinct byte hashes; an adjacent-index projection comparison has 279/437 exact IoU ties. The convention and projection checks do not independently validate physical annotation association or continuous-time recovery.

An instrumented BundleSDF replay matches normalized Point2Pose driver inputs on all 441 frames and passes 319/426/438 at 5/10/20 mm; an earlier run gives 302/425/438. Both report two tracker failures. Status precedes neural synchronization and internal processing differs; the repaired local runtime does not establish a fair ranking or physical surface accuracy.

A post-hoc $200\,\mathrm{ms}$ selected-frame gap split gives 35 gap-following and 405 other frames. At 5 mm, Point2Pose/BundleTrack pass 22/3 and 261/46 frames in the two groups. This selected grouping does not establish a causal gap effect or physical occlusion recovery.

On the same bleach-hard frames, corresponding-vertex ADD at 10 mm passes 195/45 frames for Point2Pose/BundleTrack, reversing the 10 mm ADD-S count ordering. On 438 mutually available frames, annotation-gauge rotation differences have medians $5.33^\circ/12.82^\circ$, while object-origin translation differences have $p_{95}$ 21.28/17.60 mm. ADD can penalize possible symmetries; none of these annotation-dependent diagnostics establishes physical task error or causation.

In a post-hoc BundleTrack depth-filter intervention on bleach-hard, the original/repeated 95th-percentile runs pass 49/72 frames at 5 mm, while two runs with the filter disabled pass 53/54. At 10 mm the counts are 418/418 and 415/416. These stateful runs are not resource-controlled and do not isolate the cause of the difference from Point2Pose or establish a production improvement.

A same-host, same-seed Point2Pose repeat on bleach-hard produces byte-identical saved boundaries, poses and loss flags across all 441 frames. This one repeat does not establish cross-device determinism or seed robustness.

A post-construction NumPy seed-42 intervention changes 419/441 Point2Pose pose matrices on bleach-hard, with 268/420 passes at 5/10 mm. A RANSAC-choice wrapper reproduces that output at seed 42 and the original at seed 0. Process-global replacement and divergent histories limit attribution; these post-hoc runs do not establish a robust ranking.

Among sugar's 10 mm ADD-S passes, 522/617 Point2Pose and 507/794 BundleTrack frames have annotation-gauge rotation difference above a post-hoc $5^\circ$. On tomato, the longest CAD bounding-box axis is local Z: among ADD-S passes, its directed angle $p_{95}$ is 15.84/21.05 degrees for Point2Pose/BundleTrack, compared with full-rotation $p_{95}$ 28.77/38.28 degrees. These metrics depend on initial gauge alignment and annotations, do not establish object symmetries or a physical pouring feature, and do not measure independent physical error. No selected result establishes a dataset-wide or end-to-end ranking.

\begin{table}[t]
\centering
\caption{Selected reference-mask development runs. $n$ is the post-initial denominator; P2P and BT are Point2Pose and BundleTrack. Counts include unavailable frames as failures; AUC is normalized over 0--100 mm. The 10 mm threshold is post-hoc. *Bleach-hard was selected subsequently and its sorted-file reference association is unverified.}
\label{tab:pose}
\input{pose-table.tex}
\end{table}

A CAD-assisted FoundationPose~\cite{foundationpose} diagnostic yields 934/1,001 raw sugar matrices within 10 mm. Its extra model input and missing loss flag preclude a matched comparison.

\subsection{Shared generated-mask diagnostic}
In a separate offline condition, Point2Pose receives a reference-derived initialization on frame 0 and generates subsequent masks. BundleTrack then consumes the saved Point2Pose-generated masks on the same RGB-D sequence. This is a shared segmentation source rather than a comparison of independent segmentation systems; Point2Pose's internal mask boundary was not independently captured. BundleTrack receives nonempty saved masks even when Point2Pose reports loss, while empty masks are counted as unavailable. Table~\ref{tab:generated} includes all post-initial frames. These runs use separate temporal histories from the reference-mask condition, so their differences do not isolate a causal mask effect.

At 10 mm BundleTrack has more passing frames on each of these five selected development sequences; this does not establish a general ranking or compare complete systems. On bleach, the generated-mask BundleTrack run passes 644/662 versus 301/662 in the separate reference-mask run. On sugar, the generated-mask counts at 10 mm are 332/588 for Point2Pose/BundleTrack, while at 20 mm they are 998/999. These thresholds were selected post-hoc, and the supplied reference masks and pose labels are not independent physical truth.

\begin{table}[t]
\centering
\caption{Shared Point2Pose-generated-mask development runs. $n$ excludes initialization; 10 and 20 are post-hoc ADD-S thresholds in mm, with reported failures included. The methods share the saved mask source but differ in internal handling.}
\label{tab:generated}
\input{generated-mask-table.tex}
\end{table}

On an inspected moving-wrist-camera cup trace, Point2Pose reports loss on all 294 post-initial input ticks; a BundleTrack shadow returns 264 candidates. The initial eroded mask has 1,616 positive-depth pixels out of 58,839. A post-hoc mask-restricted depth-fill shadow returns 98 non-lost outputs over 295 input ticks, but inferred depth does not establish correct recovery. The moving camera violates Point2Pose's fixed-camera assumption; no independent pose truth or cutover gate is available.

On one selected YCB-M~\cite{ycbm} static-object moving-camera trajectory, both arms use per-frame reference masks and match normalized tracker-entry receipts on all 171 frames. With frame-0 initialization, the eroded bowl mask contains 67 positive-depth pixels; Point2Pose reports lost on all 170 later frames, and both methods pass 0/170 at post-hoc 10 or 20 mm annotation-relative object-origin thresholds. BundleTrack reports 155 candidates, with 103.43 mm candidate-only median error. A post-hoc frame-1 restart, where the initialization mask has 3,887 positive-depth pixels, yields Point2Pose/BundleTrack availability 166/156 out of 169 later frames, 10 mm passes 43/119 and 20 mm passes 142/135. A further post-hoc single-object wood-block replay on the same camera trajectory matches 170 tracker-entry receipts; Point2Pose/BundleTrack have 117/106 passes at 10 mm and 151/149 at 20 mm over 169 annotated frames, reversing the bowl 10 mm ordering. A second YCB-M trajectory selects a cracker object before inspecting its RGB-D or tracker outputs, reusing the post-hoc frame-1 start. Point2Pose/BundleTrack pass 154/157 frames at 10 mm and 157/159 at 20 mm, out of 169, with 168/159 available and 170 matching tracker entries. At 10 mm, corresponding-vertex ADD / nearest-vertex ADD-S pass 4/132 versus 47/136 for Point2Pose/BundleTrack on the bowl and 111/151 versus 104/149 on the wood block, each out of 169. The nearest-vertex metric is not an explicit symmetry-group test. These correlated development outcomes depend on threshold, metric, object, scene and initialization. Different starting observations, object gauges and histories change together; depth support alone is not established as the cause. Camera motion violates Point2Pose's fixed-camera assumption, timestamps are synthetic and labels depend on image/depth fitting. This is not independent physical truth, autonomous segmentation or a live cutover test. The direct backend replay bypasses the live metric-anchor gate. In an offline helper-only check on YCB-M reference masks, the existing large-and-flat rejection is triggered on 80/170 wood-block and 96/170 second-scene cracker post-initial frames; the check is not a SAM2 or complete live initialization replay and cannot establish a false-positive rate. A post-hoc shadow cutoff raises wood/cracker helper acceptance from 90/74 to 170/157 of 170 frames while retaining the observed far-depth blockers in a small historical floor stub. One in-band-only floor-labelled stub already passes both existing and shadow rules; these fitted examples do not establish floor safety or justify a live change. Across 47 source-color-matched stride-2 historical stub frames, both rules make identical decisions; independent floor labels and depth synchronization are unavailable. A remapped held-mask file has a timestamp 401.60 s from its paired source-color record, although its pixels exactly match the source mask with matching color metadata; depth synchronization remains independently unverified. In a post-hoc offline restart at the first helper-eligible wood frame, BundleTrack passes 80/126 versus 109/126 common later frames at 10 mm; Point2Pose passes 98/126 versus 107/126 at 10 mm and 116/126 versus 108/126 at 20 mm. Initialization evidence, first gauge and state history all change, so this comparison does not establish a causal delay effect. On the full 169-frame post-frame-1 window, scored availability changes from 166/169 to 126/169 for Point2Pose and from 154/169 to 115/169 for BundleTrack; the later-start arm has 42 pre-initialization frames and one reference-aligned initialization frame without an independent score. This is not a live duty-cycle or measured latency.

An 80-frame fault diagnostic zeros the supplied mask for frames 30--39 while retaining RGB-D. Point2Pose flags those frames lost but returns a held matrix; our offline shadow adapter suppresses all ten. When masks resume, its matrices agree with the control to within $4.22\times10^{-15}$ per element through frame 79. The object remained visible, so this tests missing-mask handling, not physical occlusion recovery or physical accuracy.

On a selected 552-frame bimanual chest recording with two automatic masks per frame, Point2Pose reports 550/552 and 552/552 non-lost outputs in two-object mode, versus 552/552 each in single-object runs. The two-object run returns stale matrices on lost frames 317--318; same-host repeats yield identical saved outputs. On mutually non-lost frames, the transformed initial measured mean point differs between modes by $3.36$ and $0.70\,\mathrm{mm}$ $p_{95}$ for the two objects. These are output disagreements, not physical errors or recovery evidence. For object 0, an offline BundleTrack replay and the two-object Point2Pose arm use the same recorded RGB-D and automatic mask source. Across 549 mutually available frames, initial-point disagreement has $p_{95}$ 9.82 mm and relative rotation disagreement $p_{95}$ $14.54^\circ$; a second BundleTrack replay yields 10.07 mm and $14.75^\circ$. No independent pose truth identifies which trajectory is more accurate. For object 1, two BundleTrack repeats versus Point2Pose have initial-point disagreement $p_{95}$ 1.95 mm each; the BundleTrack repeat-to-repeat $p_{95}$ is 0.14 mm. Agreement is not physical accuracy. Rotation jumps peak at $8.51^\circ$ for BundleTrack and $14.30^\circ$ for Point2Pose; both maxima bridge reported unavailable ticks. These camera-local jumps include possible camera and object motion.

A two-object replay seeds SAM~2 on frame 0. Both objects report non-lost on all 552 frames, but object 0 has IoU below 0.5 against another automatic mask for frames 306--468; one mask covers an apparent gripper. Neither mask is independent truth. A repeated run with the previously missing SAM~2 extension retains this interval and identical poses and flags, although 71 masks change; this does not establish general extension independence. In the post-hoc low-IoU interval, object 0's median initial-point and rotation disagreements between the two runs are 4.21 mm and $4.28^\circ$ on 161 mutually non-lost frames; before that interval they are 0.33 mm and $0.64^\circ$ on 306 frames. Different evolving histories prevent a causal attribution or a physical-accuracy claim.

An instrumented replay integrates object 0's generated frame-310 mask into TSDF. It contains 10,192 positive-depth pixels outside the automatic comparison mask; volume filtering means these are not counts of erroneous mesh points. Between frames 300 and 350, sampled seed-to-control mesh distance $p_{95}$ changes from 3.66 to 17.08 mm. Different histories preclude attribution to frame 310 or identification of a correct surface. At fixed pre-integration volume, pose, RGB-D and generated dense points, replacing only the frame-310 mask changes 481 TSDF values; the generated mask updates 796 additional voxels. The generated clone matches actual integration. Neither mask is physical truth, and this does not replay an alternative trajectory.

\subsection{Persistent surfaces and selection}
On bleach-hard, Point2Pose and BundleSDF online meshes at frame 438 are aligned through initial gauges to the annotated CAD, without ICP. Their normalized driver inputs match, although internal processing differs. At 10 mm, reconstruction-to-CAD / CAD-to-reconstruction sampled fractions are 62.70\%/74.88\% and 83.88\%/59.68\%, respectively. The CAD is not an independent physical scan; neither direction establishes a general ranking.

A simultaneous two-object Point2Pose replay on the first YCB-M trajectory matches both single-object arms at all 170 actual tracker entries, per object. On 169 post-initial frames, simultaneous/single Point2Pose and single-object BundleTrack have 10 mm object-origin passes 79/43/119 for bowl and 106/117/106 for wood; at 20 mm they have 135/142/135 and 153/151/149. Reference masks are supplied per frame, object identity is not independently checked, and the moving camera violates Point2Pose assumptions. The shared multi-object pipeline changes internal state and does not establish causation, deployment latency or physical accuracy. In a post-hoc geometric assignment diagnostic, none of 163 jointly available frames prefers crossed object origins, while six frames have at least one unavailable output; reference origins are separated by at least 376 mm, and the per-frame masks carry identity. Both objects meet the 10/20 mm origin thresholds on 55/126 of all 169 post-initial frames. This does not test identity recovery after occlusion.

On the separate YCB-M cracker trajectory, a new surface replay matches 170 normalized tracker-entry receipts. At the selected final frame, Point2Pose reports non-lost and BundleTrack reports candidate with current surface integration. On a 4 mm common grid, reconstruction-to-CAD / CAD-to-reconstruction fractions within 10 mm are 69.28\% / 64.34\% for Point2Pose and 97.90\% / 55.60\% for the existing causal surface. The latter uses the raw reference mask while tracker inputs are eroded. On 160 owner-reported integrated frames, a current-helper diagnostic samples 198,617 eroded-core and 4,059 raw-only fringe points; 99.60\% and 95.49\%, respectively, lie within 10 mm of annotation-aligned CAD. The mask and annotations partly derive from RGB-D/model fitting, and current helper bytes were not independently bound during the historical replay. This is an input diagnostic, not a fused-surface counterfactual. Mesh sampling, representation, annotation and visibility differences preclude a general ranking or independent physical-accuracy claim.

A separate existing BundleTrack plus causal-surface replay matches all normalized tracker-entry receipts, while the surface uses the raw registered mask. At frame 438 on a 4 mm grid, reconstruction-to-CAD / CAD-to-reconstruction fractions are 88.61\%/60.59\% for the existing surface, 62.93\%/74.01\% for Point2Pose, and 83.65\%/58.74\% for BundleSDF. A common grid does not equalize sampled meshes and supported voxel centroids. Initial-gauge-adjusted translation disagreement from the earlier BundleTrack run has $p_{95}$ 4.00 mm, not a physical error.

Reconstructing the existing 1,500-point snapshot priority reduces that surface's CAD-to-surface fraction from 60.59\% to 36.57\%; a post-hoc farthest-point selection gives 59.73\%. A support-constrained 12 mm coarse-cell selection gives 53.57\%, with selection-only p50/p95 1.479/1.490 ms versus 0.391/0.397 ms for existing priority over 30 interleaved trials. A repeat over 14 selected snapshots cannot meet the coarse candidate's support floor at frames 0, 20 and 50; its selected-voxel Jaccard is 0.9947--1.0000 over five final adjacent transitions. A shadow-only adapter returns the existing owner order on these three cases; its validated selection costs p50/p95 2.57/3.38 ms on one snapshot. This uses no physical truth and does not establish general stability, live latency or deployment readiness.

Across 12 additional selected development snapshots from three objects and two cracker mask conditions, the fraction of stored points within 5 mm of the selected subset increases on seven snapshots, decreases on two, and is unchanged under fallback on three. The bleach final snapshot has an unavailable pose; these comparisons include selected points at zero distance and measure coverage of the stored map, including possible noise, not physical accuracy or task-feature retention. These correlated samples do not establish a general improvement.

For the final cracker snapshot on a 4 mm grid, reconstruction-to-model / model-to-reconstruction fractions within 10 mm are 63.62\% / 65.96\% for Point2Pose versus 83.22\% / 55.26\% for BundleTrack. Limiting the latter to its 1,500-point budget reduces model-to-reconstruction coverage to 44.03\%. Sampling, representation and model reference affect these results; the model is not an independent scan. On the 552-frame bimanual chest recording, a private BundleTrack plus existing causal-surface replay integrates 551 observations. At frame 551, on a 4 mm initial-camera grid, BundleTrack-surface to Point2Pose-mesh and reverse nearest-centroid distances have $p_{95}$ 38.73 and 18.05 mm. The existing surface uses the registered raw mask, while Point2Pose uses an eroded mask; representation, sampling, and pose histories differ. Neither is physical surface truth. For object 1, the final directed surface distances are 25.52/19.75 mm $p_{95}$, despite close pose agreement; BundleTrack integrated 552 observations and Point2Pose one keyframe. Instrumenting Point2Pose reveals that all 551 subsequent object-1 criterion calls retained at least 26 correspondences and stayed below the configured $15^\circ$ view-direction threshold; no further keyframe was requested. A post-hoc object-1-only $3^\circ$ threshold requests keyframe 270, integrated during frame 273; its captured update changes 3,912 TSDF values. Shared tracking changes object 0's trajectory and lost flags. No independent surface truth or complete latency establishes an improvement. On the final snapshot, directed $p_{95}$ distances to the same existing-owner surface are 25.52/19.75 mm originally versus 25.61/20.47 mm after the intervention; that surface is not independent truth. Projecting five earlier meshes to later input frames gives lower shadow median depth residual on two common-pixel subsets and higher on three; target poses use those frames and are not independent truth. For this snapshot, changing the common grid to 8 mm changes the first $p_{95}$ to 56.25--56.56 mm across two mesh sample counts and seeds, illustrating weighting sensitivity. Projected to the same frame, both surfaces cover only 624 common positive-depth automatic-mask pixels; their median absolute depth residuals there are 4.45 and 10.26 mm, respectively. These are input-consistency checks, affected by occlusion and sampling, not independent surface accuracy. Five earlier snapshots projected 30 frames forward share only 207--631 valid automatic-mask pixels per target; each uses its own target-frame estimated pose. These are limited input-consistency checks rather than a blinded future-frame test.

\subsection{Depth filter and repeatability}
On sugar, repeating the BundleTrack filter configuration yields 170 corresponding-vertex ADD passes at 10 mm, versus 174 originally; two filter-disabled runs give 212 and 192. A same-host Point2Pose repeat and a run changing the configured seed from 7 to 42 both yield byte-identical saved outputs. However, the RANSAC constructor resets the global NumPy seed to zero. A separate post-construction global NumPy seed-42 intervention yields 192 ADD and 562 ADD-S passes, with one lost frame. A private wrapper using a dedicated seed-42 stream during RANSAC choices yields byte-identical outputs to that global intervention; the same wrapper with seed zero reproduces the unwrapped original. The global state and subsequent call patterns may affect other components. These controls cover one inspected run and environment, not general seed robustness or cross-device determinism.

On the 552-frame bimanual chest replay, all four private runs report one non-candidate. Across 551 mutually available frames, initial-gauge translation-difference $p_{95}$ is 3.37 and 4.68 mm between same-setting repeats, versus 9.93--13.31 mm across settings. These differences are not physical errors.

On sugar, filtered runs achieve 794/1,001 and 729/1,001 passes at 10 mm; no-filter runs achieve 737/1,001 and 417/1,001. In frames 401--900, the no-filter runs differ by 236 passing frames with a median ADD-S difference of 1.05 mm. Correlated frames near the threshold can amplify count changes. Neither variation nor filter benefit is causally established.

\section{Limitations and Next Evaluation}
Six single-object public recordings, two additional moving-camera sequences with reference masks and one mask-assisted two-object bimanual replay cannot validate general multi-object recovery, mobile-camera geometry, independent segmentation, complete surfaces, or task success. The bimanual replay lacks independent physical pose and surface truth; its filter study uses an offline owner rather than the complete live path. Two repeats per setting do not characterize run-to-run variability. Matched full-session latency and resource limits remain unmeasured. On 552 chest frames, two offline worker repeats including IPC have $p_{50}/p_{95}$ 142.97/175.96 and 142.25/176.26 ms for two-object Point2Pose, versus 190.91/214.13 and 192.57/214.57 ms for serial two-worker BundleTrack with surface updates. A separate parallel dispatch of the same two BundleTrack/surface workers gives 134.94/158.63 and 134.50/155.59 ms, with equal candidate and integration counts. The legacy registry implementation serializes one native session per process, so this two-worker arrangement is a research setup rather than a deployed multi-object schedule. The current worktree hard-disables that registry in the production route; in a 40-frame opened regression replay it reported disabled on every frame and produced no native trace. Project instructions still identify it as the active owner and require a scored cutover, so this discrepancy is unresolved and no replacement is claimed authorized. At one additional final snapshot, serial-to-parallel and reverse supported-surface distances have $p_{95}$ 1.79/1.81 mm for object 0 and 0.92/0.93 mm for object 1; no independent surface truth is available. Environments, contention and surface work differ; none measures the live schedule or establishes a cutover. A further instrumented offline replay matches normalized RGB, metric depth, object-0 mask and intrinsics at both methods' actual tracker entries on all 552 frames. On 549 mutually available frames, their transformations of one initial masked-depth mean differ by 1.42/10.12/12.72 mm median/p95/max and their rotations by 1.64/14.68/17.19 degrees. This is inter-method disagreement with automatic masks, not independently measured physical error. A score-only adapter now converts existing registered object inputs to Point2Pose arrays; reconstructed development inputs match the saved two-object tracker-entry receipts on all 552 frames. A subsequent offline adapter-driven Point2Pose replay matches all 552 direct-input tracker entries, loss flags and returned poses. The adapter does not itself run the live session, verify identity or publish poses. An additional offline single-object replay initializes the existing SAM2 video-memory session from one automatic mask and propagates it over 552 frames. Both pose trackers receive matching actual entry arrays on all frames; BundleTrack reports 551 candidates and Point2Pose 552 non-lost outputs. Their initial masked-depth mean differs after transformation by 1.56/7.04/10.36 mm median/p95/max over 551 mutual outputs. The full disabled registry and its initialization gate are not replayed, and no independent physical truth is available. With the same saved SAM2 output, new causal-surface and Point2Pose TSDF runs have 552 matching tracker entries. At the final frame, 4 mm-grid nearest-centroid fractions within 10 mm are 80.51\% from the supported point map to the mesh and 70.13\% in reverse. The representations and fusion masks differ; neither surface is ground truth. A post-hoc head-camera depth check at that frame uses an unpromoted cross-camera transform and 5,703 automatic-mask pixels. With zero/one-pixel raster radii, the common hit count is 1,036/5,321 and the map/mesh common-hit depth residual p95 is 44.73/26.46 mm and 12.85/11.18 mm, respectively. Changing radius changes the hit set and may cross boundaries; this is not independent physical accuracy or a general ranking. In a further same-input replay, a shadow 1,500-point selector covers 67.27\% of the existing map within 5 mm versus 56.32\% for the actual owner selection. Yet on 5,703 valid head-mask pixels, their 10 mm hit counts are 495/513 at zero-pixel radius and 5,071/5,038 at two pixels (shadow/owner). This single development view does not establish a safer or more accurate selection. Across six correlated snapshots in the same recording, shadow selection raises same-map 5 mm coverage in every snapshot, while its head-depth 10 mm hit count at one-pixel radius increases once and decreases five times; at two pixels it increases five times and decreases once. Neighborhood effects and unverified cross-camera alignment prevent a general ranking. The recorded-system checkpoint reports blocked task relations and a failed latency gate.

A replacement decision requires frozen releases, reserved held-out RGB-D with verified source and mask bindings, independent physical pose/surface truth, failure-inclusive accuracy and freshness, and paired latency and jump measurements. The shadow scoreboard blocks promotion without independently verified freshness, holdout, baseline, session timing, object identity and native mask consumption. It also checks declared mask hashes on available ticks: a historical 47-frame registered-sequence replay lacks these hashes for all four arms. Hash equality alone would not verify tracker consumption. In a separate recorded 40-frame native-tap replay, every observed BundleTrack pre-frame mask matches a post-hoc 3-by-3 erosion of the saved registered mask. This capture starts after registration and does not independently bind the original SAM2 output or historical native binary. Capture also remains blocked until device-to-stream provenance is verified: matching ROS topics cannot distinguish a bag player from a sensor.

\section{Conclusion}
These selected experiments show metric- and run-dependent tradeoffs, not a best 3D reconstruction system. The evidence supports retaining the live authority while preparing an independently referenced evaluation.

\begin{thebibliography}{9}
\bibitem{bundletrack} B.~Wen and K.~Bekris, ``BundleTrack: 6D pose tracking for novel objects without instance or category-level 3D models,'' IROS, 2021. \url{https://arxiv.org/abs/2108.00516}
\bibitem{point2pose} T.-Y. Lin \emph{et al.}, ``Point2Pose: Occlusion-recovering 6D pose tracking and 3D reconstruction for multiple unknown objects via 2D point trackers,'' ECCV, 2026. \url{https://arxiv.org/abs/2604.10415v2}
\bibitem{bundlesdf} B.~Wen \emph{et al.}, ``BundleSDF: Neural 6-DoF tracking and 3D reconstruction of unknown objects,'' CVPR, 2023. \url{https://arxiv.org/abs/2303.14158}
\bibitem{tracknet} B.~Wen, C.~Mitash, B.~Ren, and K.~E. Bekris, ``se(3)-TrackNet: Data-driven 6D pose tracking by calibrating image residuals in synthetic domains,'' IROS, 2020. \url{https://doi.org/10.1109/IROS45743.2020.9341314}
\bibitem{foundationpose} B.~Wen, W.~Yang, J.~Kautz, and S.~Birfield, ``FoundationPose: Unified 6D pose estimation and tracking of novel objects,'' CVPR, 2024. \href{https://openaccess.thecvf.com/content/CVPR2024/papers/Wen_FoundationPose_Unified_6D_Pose_Estimation_and_Tracking_of_Novel_Objects_CVPR_2024_paper.pdf}{Paper}
\bibitem{ycbm} T.~Grenzd\"orffer, M.~G\"unther, and J.~Hertzberg, ``YCB-M: A multi-camera RGB-D dataset for object recognition and 6DoF pose estimation,'' 2020. \url{https://arxiv.org/abs/2004.11657}
\end{thebibliography}
\end{document}
